\section{Panorama competitivo: liderazgo tecnológico, liderazgo del ecosistema y temas de alcance nacional}

\subsection{De la competencia entre empresas a la competencia entre ecosistemas}

Cuando Lex pregunta sobre la ``ventaja competitiva'', Jensen no deja la respuesta en los parámetros del producto, sino que vuelve a la base instalada, la velocidad de los desarrolladores y la capacidad de entregar sistemas completos. Esto muestra que define al competidor como un ``ecosistema''. Cuando la unidad de competencia pasa de la empresa al ecosistema, el eje de la estrategia se desplaza del ritmo de lanzamientos a corto plazo hacia los efectos de red a largo plazo.

\begin{knowledgebox}{Estructura de competencia en tres capas}
El panorama actual puede entenderse en tres capas:
\begin{itemize}
    \item Primera capa: desempeño de chips y sistemas (capa de producto).
    \item Segunda capa: desarrolladores y pila de software (capa de plataforma).
    \item Tercera capa: cadena de suministro y coordinación industrial (capa de infraestructura).
\end{itemize}
De estas tres capas, las dos últimas determinan en mayor medida la participación de mercado a largo plazo.
\end{knowledgebox}

\subsection{Liderazgo tecnológico y el tema de la seguridad nacional}

En la parte final de la entrevista aparece la discusión sobre ``technology leadership'' y ``national security''. Lo relevante aquí no es la consigna, sino el hecho: las plataformas de cómputo ya son infraestructura crítica que incide en la investigación científica, la industria, las finanzas, la defensa y los servicios públicos. Por eso las empresas de plataforma asumen a la vez una responsabilidad comercial y una responsabilidad pública.

\begin{warningbox}{Riesgo narrativo: reducir la competencia tecnológica a un juego de suma cero}
El liderazgo tecnológico sí involucra intereses nacionales, pero si el problema se reduce a un enfrentamiento de una sola dimensión, se pasa por alto la interdependencia real de las cadenas de suministro globales. La vía más viable suele ser: mantener la cooperación dentro de los límites de la seguridad y mantener la competencia dentro de los límites de la cooperación.
\end{warningbox}

\subsection{Resumen de la sección}

La forma en que NVIDIA compite ya se extendió de la ``competencia de producto'' a la ``competencia de ecosistema e infraestructura''. Esto también explica por qué la dirección de la empresa, en su narrativa pública, habla a la vez de tecnología, industria y política.

